\chapter{CONCLUSION}\label{chapter-vi.-conclusion}

\section{Summary of Findings}\label{summary-of-findings}

This thesis investigated whether contextual bandit algorithms can replace static underwriting rules for health insurance in Cambodia, and whether Population Stability Index guardrails can ensure the resulting portfolio remains actuarially sound. Four experiments, conducted on a synthetic 2,000-applicant Cambodia dataset (20 independent seeds, seeds 1--20, N = 5,000 rounds per seed; trajectory figures use SEED = 42 unless noted), yield four principal findings.

\textbf{Finding 1 --- Bandits substantially outperform static underwriting (EXP-005).} Averaged across 20 independent seeds, LinUCB accumulates \$90,540 in cumulative reward over 5,000 rounds versus \$72,292 for the Static XGBoost baseline, a \textbf{25\% improvement} (Wilcoxon p \textless{} 0.001, Cohen's d = 2.98, large effect). Late-round average regret falls to \$2.20/round versus \$9.01/round for the baseline (Wilcoxon p \textless{} 0.001), confirming that the performance gap compounds rather than narrows as the bandit converges. (\emph{Ceiling note:} AlwaysRATED, an inadmissible constant, leads all learners by a further 30--33\%; §6.2.)

\textbf{Finding 2 --- Adaptive underwriting does not introduce demographic bias (EXP-006).} Without any explicit fairness constraint in the reward function, LinUCB achieves regional approval rates with a min/max parity of \textbf{85.72\%} across all eight regions and \textbf{90.12\%} across all seven occupations, both satisfying the EEOC four-fifths (≥ 80\%) rule (20-seed mean, converged phase rounds 3,000--4,999). Maximum 500-round sliding-window PSI remains GREEN for region (0.082 mean, 95\% CI {[}0.073, 0.091{]}) and within the low AMBER band for occupation (0.123 mean, 95\% CI {[}0.095, 0.161{]}), confirming that the approved portfolio's demographic distribution does not deviate materially from the applicant pool at any point during the run. (\emph{Scope note:} EXP-006 audits LinUCB's own parity; no parity comparison against the Static XGB baseline is computed.)

\textbf{Finding 3 --- Online adaptive linear estimation, whichever exploration wrapper is used, decisively beats the non-adaptive Static XGB baseline (EXP-007, EXP-011).} Averaged across 20 seeds, LinTS achieves a cumulative regret of \$21,149 and LinUCB \$22,774; the difference between the two uncertainty-directed bandits is not significant under a paired Wilcoxon test (p = 0.87, Cohen's d = 0.26, small effect). Both are dramatically lower than Epsilon-Greedy (\$38,281) and Static XGB (\$42,548), with pairwise differences highly significant after Bonferroni correction (p \textless{} 0.001, \textbar d\textbar{} \textgreater{} 2.7 in all cases). The follow-up ablation (EXP-011, §5.7) refines the causal attribution: a Greedy-only variant of LinUCB (α = 0, no exploration bonus) attains \$91,261 cumulative reward --- statistically indistinguishable from full LinUCB at α = 1.0 (\$90,540; paired Wilcoxon p = 0.58, d = 0.13) --- yet still beats Static XGB by +\$18,969 (p \textless{} 0.001, d = 3.01). The load-bearing mechanism on this dataset is therefore \emph{online ridge-regression updating of per-action coefficients}; UCB exploration adds no statistically significant gain on top. LinTS's operational advantage over LinUCB is parameter-free deployment, not statistical superiority. (\emph{Ceiling note:} all learners remain below AlwaysRATED under every reward model and drift scenario tested; §6.2, §5.11.)

\textbf{Finding 4 --- Human-in-the-loop underwriting improves reward at low cost (EXP-008).} Replacing the REFER arm's mathematical oracle with a simulated human underwriter (conservatism c = 0.7) increases cumulative reward by \textbf{14.8\%} (20-seed mean: \$103,951 ± 1,852 vs \$90,540 ± 5,382; Wilcoxon p \textless{} 0.001, d = 2.65) while incurring a mean review cost of \$2,263 (2.2\% of cumulative reward, well below the 15\% threshold). Human review was required on only 65 of 5,000 rounds (1.3\%) on average across 20 seeds, demonstrating that a contextual bandit naturally learns when to refer and when to decide autonomously. (\emph{Caveat:} the HITL-vs-AlwaysRATED rebenchmark rests on SEED = 42 and on EXP-008's deterministic reward basis, which is not comparable to the 20-seed stochastic mean above; on that basis HITL narrows the gap to the constant to \textless{}1\% but does not overturn it; §5.4.2.)

Taken together, these findings demonstrate that contextual bandits outperform every \emph{admissible} alternative --- decisively so against the static and uniform-exploration baselines --- and that a constant-policy ceiling bounds but does not invalidate this contribution. The deployment recommendation is therefore \textbf{conditional}: the result is best read as a rigorous proof of algorithmic feasibility --- a frozen underwriting rule leaves money on the table, and this thesis provides the machinery to measure exactly how much --- rather than as an operational prescription to deploy a contextual bandit.

\section{Limitations}\label{limitations}

Several limitations bound the generalisability of these findings.

\textbf{Synthetic data.} All 2,000 applicants were generated by a stochastic simulator calibrated to Cambodia-specific demographic and disease prevalence estimates (Hepatitis B, TB, occupational mortality multipliers). The simulator reproduces known statistical relationships but cannot capture the full heterogeneity of a live applicant pool, including rare comorbidities, fraud patterns, and socioeconomic shocks absent from the generative model.

\textbf{Single-period rewards.} The actuarial reward simulator computes a one-shot premium minus expected claim cost per applicant per round. In practice, health insurance is a multi-period contract: retention, claim experience, and renewal pricing create intertemporal dependencies that a single-round reward function does not model. A customer accepted at STANDARD rate in round 1 may file a large claim in round 1,200; this feedback loop is absent from the present experiments.

\textbf{Stationary environment.} EXP-005 through EXP-007 assume a stationary applicant distribution throughout the 5,000 rounds. Real insurance portfolios experience demographic drift --- seasonal disease spikes, economic migration, aging cohorts --- that can invalidate a bandit's learned coefficients. EXP-008 partially addresses this by introducing human overrides as a correction mechanism, but does not model environmental non-stationarity explicitly.

\textbf{Sample size.} 2,000 applicants is sufficient to demonstrate convergence on this synthetic problem but is small relative to the hundreds of thousands of policies held by even mid-sized emerging-market insurers. Bandit convergence rates scale with context dimensionality; the 34-feature space used here may require substantially more rounds to converge in a live setting with richer feature sets.

\textbf{Constant-policy ceiling.} A post-hoc analysis identified an inadmissible constant policy --- AlwaysRATED (rating 100 \% of applicants at the +25 \% loading) --- as a theoretical performance ceiling above both bandits: \$122,287 cumulative reward versus LinUCB's \$91,864 on the 20-seed common-random-number ladder harness (paired Wilcoxon p \textless{} 0.0001, Cohen's d = −5.1; Finding 1's \$90,540 is the same algorithm on EXP-005's own pipeline --- §5.0.1 reconciles the two figures). This ranking is invariant across six reward calibrations and the full demand-elasticity range; DiscountedLinUCB (\(\lambda \in \{0.999, 0.995\}\); EXP-015) and a human-in-the-loop wrapper do not close the gap. Rating every applicant at a flat surcharge is neither commercially viable --- it discourages low-risk uptake and invites adverse selection --- nor regulatorily admissible under standard underwriting governance. The constant therefore bounds the positive contribution of this thesis: the bandits beat every \emph{admissible} alternative (AlwaysRATED is technically deployable but not admissible; §5.0.1), and the ceiling quantifies the headroom remaining above their converged policies.

\section{Future Work}\label{future-work}

\subsection{Non-Stationary Drift and Adaptive Monitoring}\label{non-stationary-drift-and-adaptive-monitoring}

The most pressing extension is equipping the bandit with explicit mechanisms to detect and adapt to distributional drift in the applicant population. Cambodia's disease landscape is subject to predictable seasonal shocks --- monsoon-season TB spikes in highland provinces, Hepatitis B outbreaks in coastal areas following flooding --- that can shift the feature distribution and invalidate the bandit's learned reward coefficients within weeks.

A Discounted LinUCB (DiscountedLinUCB) addresses this through exponential forgetting: the sufficient statistics \(A\) and \(b\) are discounted at each round by a forgetting factor \(\lambda \in (0, 1)\),

\[A_t \leftarrow \lambda A_{t-1} + x_t x_t^\top, \quad b_t \leftarrow \lambda b_{t-1} + r_t x_t,\]

so that older observations contribute less to the current parameter estimate. Smaller \(\lambda\) produces faster adaptation at the cost of higher variance in stationary periods; \(\lambda \to 1\) recovers standard LinUCB.

Critically, drift detection must precede algorithmic adaptation: an insurer cannot trigger a model retrain without knowing that drift has occurred. Population Stability Index provides a natural detection layer --- applied here to the demographic distribution of the \emph{approved} portfolio (§4.8), it would in a production deployment also be applied to the \emph{incoming} applicant distribution to flag environmental shifts before they cascade into stale reward coefficients. EXP-009 (§5.9) provides preliminary empirical evidence for the underlying premise that adaptive linear bandits absorb a moderate Cambodia-calibrated shock (doubled TB prevalence and a 30\% drop in garment-worker income at round 1,500) substantially better than the non-adaptive Static XGB baseline: post-shock per-round regret for standard LinUCB and LinTS is approximately 0.29× their pre-shock level, while Static XGB recovers only to 0.85× --- a gap consistent with the bandit's online re-weighting of feature coefficients in response to the new applicant mix. EXP-009 does not, however, isolate adaptation from late-stage convergence (the caveat in §5.9), nor does it test the DiscountedLinUCB extension described below; both are matters for follow-on work.

A recommended production monitoring architecture would combine both signals: PSI computed on a rolling 3-month ingestion window triggers an alert when it crosses the 0.25 RED threshold; the alert would initiate a DiscountedLinUCB warm restart with a moderate forgetting factor applied to the current sufficient statistics rather than a full cold restart from \(A = I, b = 0\). This would preserve prior knowledge of stable subpopulations while accelerating adaptation to the drifted subgroup. DiscountedLinUCB was evaluated at \(\lambda \in \{0.999, 0.995\}\) in EXP-015 (§5.9.5); at neither value did it close the constant-policy gap (Table~20). It remains a candidate for production deployments facing more severe distributional shift than tested here, subject to further evaluation. A side-by-side comparison under PSI-triggered restart remains a concrete next step that builds directly on EXP-009's drift testbed and the PSI machinery already in \texttt{healthrl/}.

\subsection{Neural Bandit Extensions}\label{neural-bandit-extensions}

The linear reward model underlying LinUCB and LinTS assumes that the expected reward is a linear function of the context features. For a 34-feature dataset, this assumption is reasonable and produces strong empirical performance. However, real-world applicant data --- including telematics, claims history, and social network signals --- may exhibit non-linear feature interactions that linear models cannot capture. NeuralUCB and NeuralTS replace the linear estimator with a deep network while preserving the UCB or Thompson Sampling exploration strategy, at the cost of substantially increased computational overhead. A hybrid neural-linear architecture, in which a neural network learns a low-dimensional embedding and a linear bandit operates in the embedding space, offers a practical middle ground for deployment on the compute budgets typical of emerging-market insurers.

\subsection{Live Deployment and A/B Testing}\label{live-deployment-and-ab-testing}

Translating the present results into a live deployment requires a controlled A/B testing protocol in which the bandit policy is applied to a randomised subset of incoming applicants while the static rule applies to the remainder. This design allows regret to be measured against a known baseline in production, rather than against a synthetic oracle. The key operational challenge is delayed reward feedback: health insurance claims typically emerge 6--24 months after policy issuance, creating a reward delay that violates the standard bandit assumption of immediate feedback. Survival analysis models or proxy rewards based on early claims indicators would be required to close this feedback loop on a commercially viable timescale.

\subsection{Multi-Period Customer Value}\label{multi-period-customer-value}

The reward function used in this thesis captures single-period underwriting profit. A richer model would incorporate customer lifetime value: the expected net present value of all future premiums minus claims, conditioned on the customer remaining in the portfolio. Retention probability, cross-selling opportunity (life, motor), and churn risk are all contextual signals that could be incorporated into the bandit's feature vector without changing the algorithm architecture, extending the framework from a pure underwriting tool to a full customer relationship optimisation engine.

\section{Final Remarks}\label{final-remarks}

This thesis has demonstrated that contextual bandit algorithms --- specifically LinUCB, LinTS, and their human-in-the-loop extension --- can learn actuarially sound underwriting policies from interaction data, outperforming static rule-based baselines by 25\% in cumulative reward (Cohen's d = 2.98 across 20 seeds) while maintaining demographic parity across Cambodia's regional and occupational subgroups under both PSI and the EEOC four-fifths rule. PSI monitoring provides a lightweight, interpretable guardrail that signals when the approved portfolio has drifted too far from the reference distribution, triggering the human or algorithmic intervention required to keep the policy current.

The broader implication is structural. Traditional insurance underwriting in emerging markets relies on rules calibrated on historical loss experience from mature markets, applied to populations whose disease profiles, occupational structures, and economic trajectories differ fundamentally. A contextual bandit learns directly from the population it serves, updating continuously as the portfolio evolves. Combined with PSI guardrails and human-in-the-loop escalation for borderline cases, this architecture offers a path to underwriting that is more accurate and more adaptable than the \emph{admissible} static alternatives it replaces; the inadmissible constant-policy ceiling documented in §6.2 bounds this contribution and frames the work as a proof of algorithmic feasibility rather than a deployment prescription.

The work presented here is a proof of concept on synthetic data. Its value lies not in the specific numbers --- which will differ on real Cambodian claims data --- but in the demonstration that the full pipeline is feasible: data generation, bandit training, fairness auditing, benchmark comparison, human integration, and drift monitoring can all be operationalised within the resource and regulatory constraints of an emerging-market insurer. The next step is a pilot with a willing partner.

